\documentclass[12pt]{article}
\usepackage[utf8]{inputenc}
\usepackage[T1]{fontenc}
\usepackage{hyperref}
\title{Wariyangga Language Resources for K-12 Curriculum Development}
\author{Compiled by Claude Sonnet 5.5 (claude-sonnet-5-5) for Andr\'as Kornai}
\date{}
\begin{document}
\maketitle

\section{Introduction}
Wariyangga (ISO 639-3 \texttt{wri}, Glottocode wari1262; also spelled Warriyangga, Warriyangka, Wariangga) is one of four varieties of the Mantharta dialect cluster of Western Australia, along with Tharrkari, Thiin and Jiwarli. Glottolog lists it as extinct, citing Austin 1988 as ``dormant''. Wikipedia says it was still spoken around 1973. It should not be confused with the Wariangga entry in the South Australian Museum archive, which sources list separately. One source reports that it was among 20 languages prioritised in a First Languages Australia support project (2020). No school teaching or community-produced learning material was found.

\section{Primary Educational Resources}
None found.

\section{Linguistic and Reference Materials}
\begin{itemize}
\item Glottolog entry for Wariyangga: \url{https://glottolog.org/resource/languoid/id/wari1262}
\item Wikipedia article on Warriyangga: \url{https://en.wikipedia.org/wiki/Warriyangga}
\item Wikipedia article on the Mantharta languages: \url{https://en.wikipedia.org/wiki/Mantharta_language}
\item AIATSIS Austlang entry for the Mantharta group: \url{https://collection.aiatsis.gov.au/austlang/language/w28}
\item Peter Austin's \emph{A Reference Grammar of the Mantharta Languages} (2015) is described as unpublished. No URL for it was seen.
\end{itemize}

\section{Religious and Community Resources}
None found.

\section{Digital and Audio Resources}
None found.

\section{Curriculum Development Context}
Very little is available. The best prospects are archival and academic sources on the Mantharta cluster, particularly Austin's grammar and AIATSIS holdings, and any future community work. Any curriculum would be a heritage and awareness one, built with Aboriginal communities in the region and with the related varieties as a basis.
\end{document}